% Diplomatic French for printed p.217.
% Printed Y in (5.6.11) and singular identique are preserved here.
\clearpage
\SourcePage{217}

\LayerTarget{tohoku:proof:lemma-5.6.2}
\LayerTarget{tohoku:construction:equivariant-injective-envelope}
\noindent dans un module $M(A)$, qui est un foncteur (non additif !) de $A$,
l'homomorphisme $A\longrightarrow M_1(A)$ étant fonctoriel ; on itère le
procédé transfiniment, et on passe à la limite inductive, ce qui nous donne un
objet injectif $M(A)$ qui est encore un foncteur en $A$, et une injection
$A\longrightarrow M(A)$ qui est un homomorphisme fonctoriel. Cette
construction est bien définie, une fois choisi un générateur $U$ de $C$ et un
cardinal convenable. De plus, remplaçant au besoin $U$ par la somme directe de
ses transformés par $G$ (comparer avec la démonstration de prop. 5.1.1), on
peut supposer que $U$ provient d'un élément de $C(G)$. On voit alors, pour des
raisons évidentes de ``transport de structure'', que si $A$ est un objet de
$C(G)$, $G$ opère aussi dans $M(A)$, qui se trouve donc être un objet de
$C(G)$ en réalité, et l'injection $A\longrightarrow M(A)$ est compatible avec
les opérations de $G$. C'est l'injection cherchée de $A$ dans un objet de
$C(G)$ qui est injectif dans $C$.

\LayerTarget{tohoku:conclusion:lemma-5.6.2}
Cela achève la démonstration ; le lecteur qui conserverait des doutes
explicitera la démonstration dans le cas $C=C^O$ envisagé dans le lemme, en se
bornant s'il préfère au cas ``purement algébrique'' où $X$ est réduit à un
point.

\LayerTarget{tohoku:transition:main-equivariant-ext-result}
Nous sommes maintenant en mesure d'énoncer le résultat principal de ce
numéro, démontré dans les considérations qui précédaient :

\LayerTarget{tohoku:theorem:5.6.3}
\noindent\textsc{Théorème 5.6.3.}\quad \textit{Soit $X$ un espace muni d'un
groupe d'homéomorphismes $G$ et d'un $G$-faisceau d'anneaux $O$. Soit $A$ un
$G$-$O$-Module. On peut trouver, sur la catégorie $C^{O(G)}$ des
$G$-$O$-Modules, trois foncteurs spectraux cohomologiques, aboutissant au
foncteur gradué $\bigl(\operatorname{Ext}_{O,G}^n(X;A,B)\bigr)$, et dont les
termes initiaux sont respectivement :}

\LayerTarget{tohoku:equation:5.6.11}
\begin{equation}
\begin{aligned}
\mathrm I_2^{p,q}(B)
  &=H^p\bigl(Y,\mathbf{Ext}_{O,G}^q(A,B)\bigr),\\
\mathrm {II}_2^{p,q}(B)
  &=H^p\bigl(G,\operatorname{Ext}_O^q(X;A,B)\bigr),\\
\mathrm {III}_2^{p,q}(B)
  &=H^p\bigl(X;G,\mathbf{Ext}_O^q(A,B)\bigr).
\end{aligned}
\tag{5.6.11}
\end{equation}

\LayerTarget{tohoku:specialization:theorem-5.6.3}
Si $O=A$ les deux premières suites spectrales se réduisent aux suites
spectrales du théorème 5.2.1, la troisième suite spectrale est nouvelle. Si
$G$ se réduit à l'élément neutre, alors les suites spectrales I et III sont
identique, et identiques à la suite spectrale du théorème 4.2.1, tandis que la
suite spectrale II est triviale. Si $X$ se réduit à un point, les suites
spectrales II et III sont identiques, et d'ailleurs non triviales en général
(et, il semble, utiles), tandis que la suite spectrale I est
triviale.\textsuperscript{11)}

\LayerTarget{tohoku:consequence:edge-homomorphisms-five-term}
On tire des suites spectrales précédentes des ``edge-homomorphisms'' et des
suites exactes à 5 termes, que nous n'écrirons pas.

\LayerTarget{tohoku:transition:degeneration-case}
Nous n'expliciterons qu'un seul cas particulier de dégénérescence :

\LayerTarget{tohoku:footnote:11}
\noindent\textsuperscript{11)} La suite spectrale non triviale obtenue aboutit
à $\operatorname{Ext}_U(A,B)$ et son terme initial est
$H^p\bigl(G,\operatorname{Ext}_O^q(A,B)\bigr)$ ($U$ étant l'anneau engendré
par $\mathbf Z(G)$ et le $G$-anneau $O$, pour les relations de commutation
$g\lambda g^{-1}=\lambda^g$ pour $\lambda\in O$, $g\in G$). Signalons qu'on
peut obtenir cette suite spectrale plus facilement directement, en utilisant
des résolutions projectives de $A$ au lieu de résolutions injectives de $B$.
On voit en effet très facilement que si $A$ est un $U$-module libre, alors
$A$ est un $O$-module libre, et $\operatorname{Hom}_O(A,B)$ est un $G$-module
$\Gamma^G$-acyclique.

\TohokuContinuation{218}
